\documentclass[10pt,a4paper]{amsart}
\usepackage[margin=19mm]{geometry}
\usepackage{amsmath,amssymb,mathtools}
\usepackage[scr=rsfs,cal=euler]{mathalfa}
\usepackage{xfrac}
\usepackage[unicode,hidelinks,bookmarks=false]{hyperref}
\newcommand{\ie}{i.e.\ }
\newcommand{\cf}{cf.\ }
\newcommand{\resp}{respectively\ }
\newcommand{\Subsection}{\subsection}
\providecommand{\nobreakdash}{\nobreak\hbox{-}}
\setlength{\emergencystretch}{2em}
\multlinegap0pt
\usepackage{float}
\usepackage{amssymb}
\usepackage{mathtools}
\usepackage{enumerate}
\newtheorem{proposition}{Proposition}
\title{Arc shift move and region arc shift move for twisted knots}
\author{Tumpa Mahato, Prabhakar Madeti}
\date{Source: arXiv:2602.06432v1 (CC BY 4.0)}
\begin{document}
\maketitle

\noindent\emph{Source and licence.} Arc shift move and region arc shift move for twisted knots. Version arXiv:2602.06432v1 (CC BY 4.0), distributed by arXiv under the Creative Commons Attribution 4.0 International licence, http://creativecommons.org/licenses/by/4.0/, as recorded in the arXiv licence metadata for this exact version and shown on the abstract page https://arxiv.org/abs/2602.06432v1. Full licence notice: \texttt{licenses/arxiv-2602.06432-CC-BY-4.0.txt}.\par\smallskip

\noindent\textit{Editorial scope.} Counted unit: the article's proposition pr:for3 (source0.tex lines 677--697). The article's complete original proof of it is delivered with it (source0.tex lines 698--787, one \texttt{proof} environment of 90 lines). No mathematics is written, repaired or re-derived: the statement and the proof are the article's own lines, in the article's own notation, and no retained line is substituted or re-flowed.  No further source-local context is needed: every cross-reference of every retained line resolves inside the two delivered spans.  One declared adaptation: exactly 1 ancillary strings inside the retained spans are omitted (image-inclusion commands and, where the source repeats a label, the repeated label definitions) and no other line differs from the source; the figure environments, with their placement, captions and labels, are kept, and the package records the omitted strings in fidelity receipts bound to the affected bodies.  No retained line contains a citation, so the wrapper carries no bibliography and no bibliography entry.  The retained lines contain the article's own diagram code, which the wrapper typesets with the standard tikz packages; no image file is delivered and the package declares no asset dependency.  Editorial formatting only, in addition: an AMS \texttt{amsart} preamble that restates the article's own declarations and macros used by the retained lines, namely the environments \texttt{proposition} on separate counters, together with the standard packages those lines need.  The retained environments print in this document as Proposition 1 in order of appearance, whereas the article prints the same items with the numbering of its own full document.  The complete native source of the article as distributed in the arXiv e-print for this exact version is bundled unchanged as \texttt{sources/194096/source0.tex}.
	\begin{proposition}\label{pr:for3}
		Let $K$ and $K'$ be two twisted knots which transforms into each
		other by a single forbidden move $F_3(\text{or} F_4$).  And $\# b(K)$ denotes the number of bars on $K$. Then the following holds.\\
			$\bullet$ If $\#b(K)$ is even, then
			$$Q_{K'}(s,t)-Q_{K}(s,t)
			=(-1)^{l_{1}} m_{1}(st-1)^{2}+(-1)^{l_{2}} m_{2}(t-1)^{2}+(-1)^{l_{3}}m_{3}(s-1)^{2},
			$$
		where $m_i \in \{0,1,2\}, l_{i}\in \{-1,1\}$ and the
		coefficients satisfy one of the following mutually exclusive conditions:
		\begin{itemize}
			\item[(i)] $m_1=m_2=m_3=0$,
			\item[(ii)] $m_1=m_2=m_3=1$ and 
			( $l_1=l_2$ or $l_1=l_3$).
			\item[(iii)] $m_i\in\{0,2\}$ for all $i$, with $m_1\neq m_2=m_3$ and
			$l_2\neq l_3$.
		\end{itemize}
	$\bullet$ If $\# b(K)$ is odd, then
			$$Q_{K'}(s,t)-Q_{K}(s,t)=
			\pm m_{4}(st-1)(s-1),	$$
			for $m_{4}\in \{0,2\}$.
	\end{proposition}

	\begin{proof}
			We prove this for $F_3$ move. A similar argument applies in the case of $F_4$ move. Consider the Gauss
		diagrams $G$ and $G'$ corresponding to $K$ and $K'$ respectively (Fig.~\ref{fig:f3case}). Let us denote the  chords (and crossings) by $c_1, c_2$ and $c'_1, c'_2$ before and after the $F_3$ move  respectively.
			\begin{figure}[h!]
			\centering
			
			\caption{}
			\label{fig:f3case}
		\end{figure}
		
		We note that $F_3$ move does not change the sign of the crossings. Then, $sgn(c_1)=\varepsilon_{1}=sgn(c'_1)$ and $sgn(c_2)=\varepsilon_{2}=sgn(c'_2)$. 
		
		From Fig.~\ref{fig:f3case}, we observe that
		\[ind(c'_{1})=ind(c_{1}) + \varepsilon_{2},\]
		\[ind(c'_{2})=ind(c_{1}) - \varepsilon_{1}.\]
		Since, there is one bar between the chord ends  $\bar{c_1}$ and $\bar{c_2}$,
		\[|p^{*}(c_{i})-p^{*}(c'_{i})|=1 \; \text{for} \; i=1,2 , \; * \in \{O,U\}.\]
		Consequently, the indices change from odd(even) to even(odd) after the move and the values of $p^*, *\in \{O,U\}$ switch between $0$ and $1$. Also,  We discuss
		all the cases based on the indices of $c_1,c_2$ of being even or odd.\\
		\noindent
		Case I: When both $ind(c_{1})$ and $ind(c_{2})$ are even.
		So, $ind(c'_{1})$ and $ind(c'_{2})$ both are odd and the corresponding values of $\bar{\rho}^{O}$ are as follows.
		\[\bar{\rho}^{O}(c_{1})=0=\bar{\rho}^{O}(c_{2}) \; \text{and} \; \bar{\rho}^{O}(c'_{1})=1=\bar{\rho}^{O}(c'_{2}). \]
		
		Therefore, \begin{align*}
			Q_{K'}(s,t)-Q_{K}(s,t) &=\varepsilon_{1}(st^{p^{O}(c'_{1})}-1)(st^{p^{U}(c'_{1})}-1)+\varepsilon_{2}(st^{p^{O}(c'_{2})}-1)(st^{p^{U}(c'_{2})}-1)\\ &-\varepsilon_{1}(t^{p^{O}(c_{1})}-1)(t^{p^{U}(c_{1})}-1)-\varepsilon_{2}(t^{p^{O}(c_{2})}-1)(t^{p^{U}(c_{2})}-1).
		\end{align*}
	 When $\# b(K)$ is even, \[p^{O}(c_i)=p^{U}(c_i)\; , i=1,2.\]
		Then for all the possible values of $\varepsilon_{1}, \varepsilon_{2}$, $p^{O}(c_i)$ and  $p^{U}(c_i)$ we get the following values.\\
		For $\varepsilon_{1}=\varepsilon=\varepsilon_{2}$, where $\varepsilon \in \{1,-1\}$,
		$$Q_{K'}(s,t)-Q_{K}(s,t)
		=\begin{cases}
			2\varepsilon(st-1)^{2}, \hspace{4cm} \text{for} \; p^{O}(c_i)=0, i=1,2,\\
			\varepsilon ((st-1)^{2}+(s-1)^{2}-(t-1)^{2}) , \quad \text{for} \; p^{O}(c_1)=0, p^{O}(c_2)=1,\\
			\varepsilon ((st-1)^{2}+(s-1)^{2}-(t-1)^{2}), \quad \text{for} \; p^{O}(c_1)=1, p^{O}(c_2)=0,\\
			2\varepsilon((s-1)^{2}-(t-1)^{2}), \hspace{2.15cm} \text{for} \; p^{O}(c_i)=1, i=1,2.
		\end{cases}
		$$
		For $\varepsilon_{1}=\varepsilon=-\varepsilon_{2}$, where $\varepsilon \in \{1,-1\}$  \\
		$$Q_{K'}(s,t)-Q_{K}(s,t)
		=\begin{cases}
	\varepsilon ((st-1)^{2}-(s-1)^{2}+(t-1)^{2}) , \quad \text{for} \; p^{O}(c_1)=0, p^{O}(c_2)=1,\\
	\varepsilon (-(st-1)^{2}+(s-1)^{2}-(t-1)^{2}), \; \text{for} \; p^{O}(c_1)=1, p^{O}(c_2)=0,\\
0, \hspace{5.7cm} \text{otherwise}.
		\end{cases}
		$$
	When $\#b(K)$ is odd, we get the following.
		For $p^{O}(c_i),p^{U}(c_i) \in \{0,1\}$,
		$$Q_{K'}(s,t)-Q_{K}(s,t)=\begin{cases}
			2\varepsilon(st-1)(s-1),  \quad \text{for} \; \varepsilon_{1}=\varepsilon=\varepsilon_{2}, \varepsilon \in \{1,-1\},\\
			0,  \hspace{3cm} \text{for} \; \varepsilon_{1}=\varepsilon=-\varepsilon_{2},   \varepsilon \in \{1,-1\}.
		\end{cases}	$$
		\noindent
		Case II:  One of $ind(c_{1})$ and $ind(c_{2})$ is odd and the other is even. Without loss of generality, let us assume that $ind(c_{1})$ is odd and $ind(c_{2})$ is even.
		So, $ind(c'_{1})$ is even and $ind(c'_{2})$ is odd and the corresponding values of $\bar{\rho}^{O}$ are as follows.
		\[\bar{\rho}^{O}(c_{1})=1=\bar{\rho}^{O}(c'_{2}) \; \text{and} \; \bar{\rho}^{O}(c'_{1})=0=\bar{\rho}^{O}(c_{2}). \]
		Therefore, \begin{align*}
			Q_{K'}(s,t)-Q_{K}(s,t) &=\varepsilon_{1}(t^{p^{O}(c'_{1})}-1)(t^{p^{U}(c'_{1})}-1)+\varepsilon_{2}(st^{p^{O}(c'_{2})}-1)(st^{p^{U}(c'_{2})}-1)\\ &-\varepsilon_{1}(st^{p^{O}(c_{1})}-1)(st^{p^{U}(c_{1})}-1)-\varepsilon_{2}(t^{p^{O}(c_{2})}-1)(t^{p^{U}(c_{2})}-1).
		\end{align*}
		When  $\# b(K)$ is even, then
		\[p^{O}(c_i)=p^{U}(c_i)\; , i=1,2.\]
		Then for all the possible values of $\varepsilon_{1}, \varepsilon_{2}$, $p^{O}(c_i)$ and  $p^{U}(c_i)$ we have the following values.\\
		For $\varepsilon_{1}=\varepsilon=\varepsilon_{2}$, where $\varepsilon \in \{1,-1\}$,
		$$Q_{K'}(s,t)-Q_{K}(s,t)
		=\begin{cases}
			\varepsilon (-(st-1)^{2}+(s-1)^{2}-(t-1)^{2}),\hspace{0.15cm} \; \text{for} \; p^{O}(c_i)=1, i=1,2,\\
			\varepsilon ((st-1)^{2}-(s-1)^{2}+(t-1)^{2}), \; \quad\; \text{for} \; p^{O}(c_i)=0, i=1,2,\\
			0, \hspace{5.95cm} \text{otherwise.}
		\end{cases}
		$$
		For $\varepsilon_{1}=\varepsilon=-\varepsilon_{2}$, where $\varepsilon \in \{1,-1\}$,
		$$Q_{K'}(s,t)-Q_{K}(s,t)
		=\begin{cases}
			\varepsilon (-(st-1)^{2}-(s-1)^{2}+
			(t-1)^{2}),  \hspace{0.5cm} \text{for} \; p^{O}(c_i)=1,  i=1,2,\\
			-2\varepsilon(st-1)^{2},  \hspace{3.85cm}\text{for} \; p^{O}(c_1)=1, p^{O}(c_2)=0,\\
			2\varepsilon((t-1)^{2}-(s-1)^{2}),  \hspace{2.4cm} \text{for} \; p^{O}(c_1)=0, p^{O}(c_2)=1,\\
			\varepsilon (-(st-1)^{2}-(s-1)^{2}+(t-1)^{2}), \hspace{0.6cm} \text{for} \; p^{O}(c_i)=0, i=1,2.
		\end{cases}
		$$
		
		
		Now, when $\# b(K)$ is odd, then for the possible values of $\varepsilon_{1}, \varepsilon_{2}$, $p^{O}(c_i)$ and  $p^{U}(c_i)$ we have following values.
		$$Q_{K'}(s,t)-Q_{K}(s,t)=\begin{cases}
			-2\varepsilon(st-1)(s-1), \quad \quad \; \text{for} \; \varepsilon_{1}=\varepsilon=-\varepsilon_{2}, \varepsilon \in \{1,-1\},\\
			  0, \hspace{3.1cm}  \quad\text{for} \; \varepsilon_{1}=\varepsilon=\varepsilon_{2}, \varepsilon \in \{1,-1\}.
		\end{cases}	$$
		
		All the remaining cases follows from the cases discussed above.
	\end{proof}

\end{document}
